\section{Background and Motivation}
\label{sec:bg_mt}
\para{Prefix caching and KV-cache tiering}
Serving engines such as SGLang~\cite{zheng2024sglang} and vLLM~\cite{kwon2023pagedattention} reuse the longest matching prefix, computing only the remaining tokens. Tiered caching systems such as CachedAttention~\cite{gao2024cachedattention}, Mooncake~\cite{qin2025mooncake}, and SGLang HiCache~\cite{sglang2026hicache} extend the KV-cache beyond device memory with host memory and optionally disk storage to improve capacity. However, off-device KV-cache must be promoted back before use.

\para{KV-cache management for LLM agents}
Agent serving systems use recurring patterns in agentic workloads to manage the KV-cache. KVFlow~\cite{pan2025kvflow} derives each agent's steps-to-execution from the workflow to guide eviction and prefetching of each agent's prompt prefix; Helium~\cite{wadlom2026helium} analyzes the workflow to order LLM calls and prewarm static prefixes; Learned Prefix Caching~\cite{yang2025lpc}, PBKV~\cite{zheng2026pbkv} and CacheScout~\cite{zhang2026cachescout} predict upcoming agent invocations with trained models to decide which caches to retain. They operate the KV-cache on the identified prefix. Some systems, such as Pythia~\cite{yu2026pythia}, furthermore prefill the predicted next invocations to proactively prepare KV-cache that does not exist yet. 



% Learned Prefix Caching~\cite{yang2025lpc} predicts from the user prompt whether a conversation continues and combines the prediction with elapsed time. KVTether~\cite{fu2026kvtether} follows the application's message operations instead: a replaced or removed message invalidates the suffix it affects, and a forked context keeps what it still needs. 

\para{Frameworks with LLM calls as typed functions}
A class of agent frameworks exposes LLM invocations through typed programming interfaces. Such as byLLM's meaning-typed programming (MTP) abstraction~\cite{dantanarayana2025mtp}, DSPy's Signatures and Modules~\cite{khattab2024dspy}, and NOOA's typed Python methods~\cite{furgale2026nooa}. The developer declares a function or signature describing the desired behavior, inputs, and return type, and invokes it like an ordinary call within the program's control flow. At each invocation, the framework runtime lowers the declaration and its argument values into final prompts

\begin{figure}[t]
    % \vspace{-8pt}
  \centering
  \definecolor{csTask}{HTML}{0072B2}
  \definecolor{csDesc}{HTML}{009E73}
  \definecolor{csVal}{HTML}{E69F00}
  \definecolor{csSchema}{HTML}{CC79A7}
  \definecolor{csScaf}{gray}{0.55}
  \newlength{\csbar}\setlength{\csbar}{1.2pt}
  \newcommand{\csfont}{\fontsize{5}{5.8}\selectfont}
  \newcommand{\ph}[1]{{\normalfont\itshape$\langle$#1$\rangle$}}
  % \csseg{color}{text}: one component, a tinted band with a bar in its color.
  \newcommand{\csrolechip}{\ifcsfirst{\fboxsep=0.6pt\colorbox{black!62}{\sffamily\fontsize{4.3}{4.3}\selectfont\color{white}\strut\csrole}}\hspace{2pt}\global\csfirstfalse\fi}
  \newcommand{\csseg}[2]{\par\noindent
    \sbox0{\parbox{\dimexpr\linewidth-\csbar-3pt\relax}{\raggedright\hyphenpenalty=10000\strut\csrolechip#2\strut}}%
    {\fboxsep=0pt\colorbox{#1!13}{{\color{#1}\vrule width\csbar height\ht0 depth\dp0}\hspace{1.6pt}\usebox0\hspace{1.4pt}}}\par}
  \newcommand{\csscaf}[1]{\csseg{csScaf}{#1}}
  % \csmsg{role}{segments}: one message as a card, its role on a tab.
  % The role sits as a chip at the start of the card's first line instead of on a tab above it.
  \newif\ifcsfirst
  \newcommand{\csrole}{}
  \newcommand{\csmsg}[2]{\par\nointerlineskip\noindent
    \renewcommand{\csrole}{#1}\csfirsttrue
    \vbox{\offinterlineskip
      \hbox{\fboxsep=0pt\fboxrule=0.4pt\fcolorbox{black!62}{white}{\parbox{\dimexpr\linewidth-2\fboxrule\relax}{\ttfamily\csfont\offinterlineskip#2}}}}%
    \par\nointerlineskip\vspace{1.4pt}}
  \newcommand{\cshead}[1]{\par\nointerlineskip\noindent{\sffamily\scriptsize\bfseries\strut #1}\,{\color{black!35}\hrulefill}\par\nointerlineskip\vspace{1.4pt}}
  \newcommand{\cskey}[2]{{\fboxsep=0pt\colorbox{#1!13}{{\color{#1}\vrule width\csbar height 1.55ex depth 0.45ex}\hspace{1.6pt}{\sffamily\fontsize{5}{6}\selectfont#2}\hspace{1.6pt}}}}
  \begin{minipage}[t]{0.488\linewidth}
    \cshead{byLLM}
    \csmsg{system}{%
      \csscaf{This is a task you must complete by returning only the output. ...}}
    \csmsg{user}{%
      \csseg{csTask}{assess\_evidence(search\_results: str, question: str, claim: str, evidences: list[Evidence]) -> Finding --- \ph{sem}}
      \csseg{csDesc}{search\_results: str ---- \ph{sem}\newline question: str ---- \ph{sem}\newline claim: str ---- \ph{sem}\newline evidences: list[Evidence] ---- \ph{sem}}
      \csseg{csVal}{search\_results = \ph{value}\newline question = \ph{value}\newline claim = \ph{value}\newline evidences = \ph{value}}
      \csseg{csSchema}{Schema requirements: \ph{Finding fields}}}
    \csmsg{response\_format}{%
      \csseg{csSchema}{\ph{Finding JSON schema}}}
    \vspace{0.8pt}
    \cshead{NVIDIA NOOA}
    \csmsg{system}{%
      \csscaf{<strategy\_prompt> You are generating structured data matching the specified return type. ... </strategy\_prompt>}}
    \csmsg{user}{%
      \csscaf{<sys tag="1"> Task(prompt='''\# Your task}
      \csseg{csTask}{\ph{docstring}}
      \csseg{csDesc}{Args:\newline \hspace*{1em}search\_results: \ph{arg doc}\newline \hspace*{1em}question: \ph{arg doc}\newline \hspace*{1em}claim: \ph{arg doc}\newline \hspace*{1em}evidences: \ph{arg doc}}
      \csseg{csVal}{\#\# Input parameters:\newline search\_results = \ph{value}\newline question = \ph{value}\newline claim = \ph{value}\newline evidences = \ph{value}}
      \csscaf{Perform the task and return the result directly.''') </sys>}}
    \csmsg{response\_format}{%
      \csseg{csSchema}{\ph{Finding JSON schema}}}
  \end{minipage}\hfill
  \begin{minipage}[t]{0.488\linewidth}
    \cshead{DSPy}
    \csmsg{system}{%
      \csseg{csDesc}{Your input fields are:\newline 1. `search\_results` (str): \ph{desc}\newline 2. `question` (str): \ph{desc}\newline 3. `claim` (str): \ph{desc}\newline 4. `evidences` (list[Evidence]): \ph{desc}\newline Your output fields are: 1. `finding` (Finding): \ph{desc}}
      \csscaf{All interactions will be structured ...\newline [[ \#\# search\_results \#\# ]] \{search\_results\}\newline ... [[ \#\# evidences \#\# ]] \{evidences\}}
      \csseg{csSchema}{[[ \#\# finding \#\# ]] \{finding\} \# note: ... JSON schema: \ph{Finding JSON schema}}
      \csscaf{[[ \#\# completed \#\# ]]}
      \csseg{csTask}{In adhering to this structure, your objective is: \ph{docstring}}}
    \csmsg{user}{%
      \csseg{csVal}{[[ \#\# search\_results \#\# ]] \ph{value}\newline [[ \#\# question \#\# ]] \ph{value}\newline [[ \#\# claim \#\# ]] \ph{value}\newline [[ \#\# evidences \#\# ]] \ph{value}}
      \csscaf{Respond with the corresponding output fields, starting with `[[ \#\# finding \#\# ]]` ...}}
    \vspace{0.6pt}
    {\raggedright\noindent\baselineskip=7.5pt
    \cskey{csScaf}{framework scaffolding}\hspace{3pt}\cskey{csTask}{task instruction}\hspace{3pt}\cskey{csDesc}{parameter descriptions}\hspace{3pt}\cskey{csVal}{argument values}\hspace{3pt}\cskey{csSchema}{output schema}\par}
    \vspace{3pt}
    {\raggedright\sffamily\fontsize{5}{6.6}\selectfont
    \begin{itemize}[nosep, leftmargin=6pt, labelsep=2pt, label=\textbullet, topsep=0pt, itemsep=0pt]
      \item \ph{sem}: byLLM \texttt{sem} string
      \item \ph{docstring}: DSPy Signature or nooa docstring
      \item \ph{desc}: DSPy field \texttt{desc}
      \item \ph{arg doc}: nooa \texttt{Args} entry in the docstring
      \item \ph{value}: runtime argument value
    \end{itemize}\par}
  \end{minipage}
  % mlsys2025.sty hard-codes 10pt above captions; trimming it lets §2.5 start on this page.
  \vspace{-6pt}
  \caption{Prompts lowered by three frameworks for the same call share five kinds of components; only their order differs.}
  \label{fig:common-structure}
    \vspace{-8pt}
\end{figure}

Although these agent frameworks differ in surface syntax, they lower LLM invocations in the same way: the intermediate representation is lowered with the runtime argument values to produce the prompt. As shown in Figure~\ref{fig:common-structure}, these frameworks emit the same five kinds of components: framework scaffolding, task instruction, parameter descriptions, argument values, and output schema. The only difference is their relative positions. 

\emph{Scheduling Opportunity:} The order of the five components is a framework convention: the runtime applies the same template to every call. Moreover, the order of argument values is just an artifact of the developer's personal preference: it follows the order of the argument list in the code. Therefore, this layout can be managed for better performance. For example, DSPy's static adapter~\cite{dspy2026chatadapter} lets users tweak the lowering process to put static prompt components at the front. NOOA~\cite{furgale2026nooa} lets developers label arguments as static prefix, append-only history or dynamic suffix so that the framework can order them for better cache reuse.

\para{Intra-prompt heterogeneity in programs}
Prompts can be decomposed into components with different sharing patterns and lifetimes, and frameworks with LLM calls as typed functions offer structural guarantees to make this information analyzable. Because components are mapped to program variables, the semantic boundaries and relationships are therefore analyzable from the program.

\emph{Scheduling Opportunity:} 
Prior serving policies fail to fully exploit the reuse in agent workloads because they ignore the component-level structure of prompts. By reasoning about the program itself, scheduling operations on specific KV-cache entries can be determined.

\para{From Deficiencies to Design}
Ignoring intra-prompt heterogeneity costs a serving system in two ways. First, shared tokens lose prefix caching: once a differing token appears ahead of them, their KV-cache can no longer be matched. Second, KV-cache retention lose precision: the system sees only a matched prefix and cannot manage entries by their individual reuse. Both costs are substantial. Among the agents we evaluated in \S\ref{sec:eval}, 51--92\% of shared tokens miss the prefix cache because a differing token precedes them, and about 20\% of served prompt tokens are used only once within a session. The layout freedom and the analyzability of intra-prompt heterogeneity address these deficiencies directly, upon which \sysname is built.


